\documentclass[15pt]{scrartcl}
\usepackage[sexy]{evan}

\newcommand{\C}{\mathbb{C}}
\begin{document}

\title{Rep Note}
\date{}
\author{Jiawen Huang}
\maketitle
\section{Assumption}
We are working on $\C$ and finite group $G$. 
\section{Irreducible and Decomposable Representation}
Counterexample: 
\[\phi: \C\to GL_2(\C)\]
\[\phi(a)=\begin{bmatrix}
1 & a \\
0 & 1
\end{bmatrix}\]
This is irreducible but not decomposable. 
\subsection{Unitary Representation}
\begin{definition}
Given a rep $\phi: G\to GL(V)$ and $V$ is an inner product space. $\phi$ is irreducible if 
\[\langle \phi_g(v), \phi_g(w)\rangle=\langle v, w\rangle\]
\end{definition}
In $\C$, the standard inner product is defined as $\langle v,w\rangle=\sum \overline{v_x}w_x$
\begin{proposition}
If rep $\phi$ is unitary, then $\phi$ irreducible if and only if $\phi$ is indecomposable. 
\end{proposition}
\begin{proof}
If $\phi$ is reducible, then let $W$ be the $G$ invariant subspace of $V$, that is $\phi_g(W)=W$ for all $g$. Consider the subspace $W^\perp$. Notice that for any $x\in W^\perp$, any $v\in W$ and any $g\in G$
\[0=\langle x, \phi_{g^{-1}}(v)\rangle=\langle \phi_{g}(x),\phi_g\phi_{g^{-1}}(v) \rangle=\langle \phi_{g}(x),v \rangle \]
Thus, for any $x\in W^\perp $, $\phi_g(x)$ is orthogonal to $v$ for all $v\in W$ and $g\in G$. Thus, $W^\perp$ is a G invariant subspace of $V$. Thus, $V=W\oplus W^\perp$ as rep. 
\end{proof}
\begin{remark}
We would like these two conecpts to be equivalent. So we want our rep to be unitary, or equivalent to a unitrary rep. 
\end{remark}
\begin{theorem}
Every representation of a finite group $G$ is equivalent to a unitary rep. From now on, we assume all the reps are unitary. 
\end{theorem}
\begin{proof}
Let $\rho : G\to GL(V)$ be a rep. It is similar to $\phi: G\to GL(\C^n)$ where $\C^n$ has standard basis. We can find an inner product of $\C^n$.  
We define a new inner product of the space. 
\[(v,w)=\sum_{g\in G}<\phi_g(v), \phi_g(w)>\]
This is well defined because $G$ is finite. And we can verify $(v,w)$ is indeed an inner product of $\C$. 
\\

It's easy to see $\phi $ is a unitary rep w.r.t. the inner product $(v,w)$. 
\end{proof}
\begin{theorem}[Maschke]
Every rep of a finite group can be decomposed into irreducible rep.
\end{theorem}
\begin{proof}
By theorem 2.4 and proposition 2.2, rep of a finite group is irreducible is equivalent of being indecomposable. 
\end{proof}

\section{Morphism of Reps}
Does any linear transformation is morphism of rep? No, but we can turn it into a morphism of rep. 
\begin{definition}
Given rep $\phi: G\to GL(V)$ and $\rho: G\to GL(W)$. For any linear transformation $T$ from $V$ to $W$, we define \[T^{\#}:=\frac{1}{|G|}\sum_{g\in G}\rho_{g^{-1}}T\phi_g\]
\end{definition}
\begin{proposition}
\begin{itemize}
\item If $T$ is a morphism of reps, then $T=T^\#$
\end{itemize}
\end{proposition}
\begin{lemma}[Schur's Lemma]
The morphism between two irreducible reps $\phi$ and $\rho$ is either 0 or invertable. Consequently: 

\begin{itemize}
\item If $\phi\not\sim\rho$, then $\Hom_G(\phi, \rho)=0$ 
\item If $\phi=\rho$, then for any $T\in \Hom_G(\phi, \rho)$, $T=\lambda I$ 
\item if $\phi=\rho$, for any linear trasnformation $L$, $L^\#=\frac{Tr(L)}{deg(\phi)}I$
\end{itemize}
\end{lemma}
\begin{corollary}
Let $G$ be an abelian group, irreducible rep of $G$ must be degree of one. 
\end{corollary}
For any $h\in G$, $\phi_h$ commutes with $\phi_g$ for any $g\in G$, then $\phi_h\in \Hom_G(\phi, \phi)$. Then use schur lemma. 
\begin{corollary}
Any matrix with finite order must be diagonalizable. 
\end{corollary}
\section{The Orthogonality Relations}
\begin{definition}
Given a finite group $G$, define 
\[L(G)=\C^G=\{f\mid f:  G\to G\}\]
This is in fact a $\C$ vector space and we can define the inner product as \[<f_1, f_2>=\frac{1}{|G|}\sum_{g\in G}f_1(g)\overline{f_2(g)}\]
Its dimension over $\C$ is the size of $G$. 
\end{definition}
\begin{definition}
Given a rep $\phi$. Define $\phi_{ij}(g)$ to be the (i,j) entry of the matrix $\phi(g)$. Notice that $\phi_{ij}\in L(G)$. 
\end{definition}
\begin{theorem}[Schur orthogonlaity relations]

Suppose $\phi$ and $\rho$ are inequivalent irreducible unitary representation. Then, 
\begin{itemize}
\item $<\phi_{ij},\rho_{k\ell}>=0$
\item $<\phi_{ij},\phi_{k\ell}>= 
\begin{cases}
    \frac 1n & \text{if } i=k, j=\ell \\
    0 & \text{else}
    \end{cases}
$
\end{itemize}

\end{theorem}

\begin{proof}
The key fact to prove this is $(E_{ik}^\#)_{jl}=<\phi_{ij},\rho_{kl}>$. We then use schur lemma. 
\end{proof}
Thus, we know that $\phi_{ij}$ form an orthonormal set of $L(G)$. Thsi implies
\[\sum_{\text{irreducible } \phi}\deg(\phi)^2\leq |G|\]
With this, we are ready to develop character theory. 
\section{Character Theory}
\begin{definition}
\[\chi_{\phi}(g)=Tr(\phi(g))=\sum_{i=1}^n\phi_{ii}(g)\]
\end{definition}
Notice that $\chi_\phi$ is in $L(G)$.
\begin{proposition}
All of them are trivial by the definition
\begin{itemize}

\item $\chi_\phi(1)=\deg(\phi)$
\item If $\phi \sim \rho$, $\chi_\phi=\chi_\rho$
\item $\chi_\phi(gxg^{-1})=\chi_\phi(x)$, $\chi_\phi$ is a class function. 
\item $\chi_{\phi\oplus\rho}=\chi_{\phi}+\chi_{\rho}$
\end{itemize}
\end{proposition}

Denote $Z(L(G))$ as the set of all class functions
\begin{proposition}
$\dim(Z(L(G)))=\#\text{conjuagcy classes of } G$
\end{proposition}
Trivial. 



\begin{theorem}
$<\chi_\phi, \chi_\rho>=0$ iff $\phi\not\sim\rho$
\\
$<\chi_\phi, \chi_\rho>=1$ iff $\phi\sim\rho$
\\

Consequently: 
\begin{enumerate}
\item Character fully characterizes the equivalent classes of reps
\item Notice that $\chi_\phi$ is a class function so the number of irreducible reps $\leq $ number of conjugacy classes of $G$. 
\end{enumerate}
\end{theorem}
Proof is easy by Theorem 4.3. 
\\

We have shown that a finite group rep has a decomposition. Is it unique? YES. 
\begin{theorem}
Every finite group representation can be uniquely decomposed into irreducible representations. 
\end{theorem}
\begin{corollary}
A rep $\phi$ is irreducible if and only if $<\chi_\phi, \chi_\phi>=1$
\end{corollary}
This gives us a way to decide whether a rep is irreducible or not. 

\begin{problem}
Find all the irreducible reps of $S_3$. 
\end{problem}
1. $\sum \chi_\phi(1)^2=\sum \deg(\phi)^2\leq 6$
\\
2. Trivial rep
\\
This is tough!

\section{Regular Rep $L$ }

This is the rep of $G$ on the group algebra $\C[G]$. 
This is nice because we can compute $\chi_L(g)=
\begin{cases}
    |G| & \text{if } g=1 \\
    0 & \text{else}
    \end{cases}$
\\

Thus, $<\chi_L,\chi_{\phi_i}>=m_i=d_i$

\begin{corollary}
\[|G|=\sum \deg(\phi)^2\]
And $\phi_{ij}^{(k)}$ form a basis of $L(G)$. 
\end{corollary}

Furthermore, we can show that $\chi$ of all irreducible reps form a basis of $Z(L(G))$ (This is a fairly complicated result-Artin Theorem). This concludes that the number of conjugacy classes $=$ the number of irreducible reps. 












\end{document}